% Transcription — fonds Alexandre Grothendieck, shelfmark 64, batch 8 (pages 141-156).
% Established from the facsimile archives/batches/64/batch-08.pdf (a copy of 64.pdf
% fetched through the project relay, no cover sheet: PDF page k = archivists' page 20(N-1)+k),
% per the skill transcribe-grothendieck. The folder is 156 pages in eight batches.
%
% Pass: Opus 5.5 (claude-opus-5-5), 2026-09-27 — first pass, unchecked against the pages by a human.
%
% Pages skipped: 141 (blank sheet), 152 (blank sheet).
% His pagination: none visible on these pages.
% What the batch is about: pp. 142-151, the complex-analytic moduli of elliptic curves:
%   lattices, Torelli structures (bases of H_1), the upper half-plane D^+ as fine moduli
%   space for strict Torelli structures, the universal family (D^+ x C)/Z^2, the action of
%   SL(2,Z) (via the involution tau_infty, formulas (1)-(22)), M_{1,1}, the "differential"
%   rigidification and PD^+. Pages 146 and 147 are bound in the reverse of the argument's
%   order. p. 153 opens with an unrelated group-extension scrap, then (as pp. 154-155)
%   normality of quadratic extensions A[S]/(S^2-a) and quadratic covers of P^1_Z, with
%   the Legendre-type example Q(T) = T(T-lambda), lambda = 2^a 3^b. p. 156 is jottings on
%   an administrative notice stamped 9 and 16 JUIN 1981.
% Notation fixed once for the batch: \mathfrak{D}^+ for the upper half-plane (his gothic D),
%   \mathfrak{X}_{\mathfrak{D}^+} and \mathfrak{X}_{P\mathfrak{D}^+} for the modular
%   families, \tau_\infty for his involution, \mathcal{T} for Teichmüller-type groups,
%   \mathfrak{m}, \mathfrak{r}(C) for maximal ideals, \mathbb{E}^1_\mathbb{Z} as written.
% \ill{}: 57. \uncertain{}: 66.

\documentclass[11pt,a4paper]{article}
\input{../preamble/grothendieck}

\folder{64}
\batch{8}
\pages{141}{156}
\dating{[à partir de 1980- 1982]}
\watermark{Édition de démonstration}
\foldertitle{Modules courbes elliptiques en caractéristiques ≠2 : notes manuscrites (s.d.)}

\begin{document}

\page{142}
$\cdot$ Une courbe elliptique $E$ sur $\mathbb{C}$ est déterminée par
\[
(1)\quad
\begin{cases}
\text{a) « réseau » } H,\ \mathbb{Z}\text{-module libre de rang } 2 \\
\quad \bigl(H = \pi_1(E, o_E) \simeq H_1(E,\mathbb{Z}) = H^1(E,\mathbb{Z})^{\vee}\bigr) \\
\text{b) Une structure complexe sur } V = H \otimes_{\mathbb{Z}} \mathbb{R}
\ (\simeq t_E, \text{ espace tangent à l'origine de } E).
\end{cases}
\]
\note{à la fin de la ligne a), un court groupe biffé (\struck{\ill{}}), et un
signe biffé entre $H \otimes_{\mathbb{Z}} \mathbb{R}$ et « $(\simeq t_E$ ».}

La donnée a) permet de reconstituer $E$ comme groupe de Lie réel par
\[
(2)\qquad E \simeq H \otimes_{\mathbb{Z}} \mathbb{R} / H
\]
et la donnée b) redonne la structure complexe sur $E$, par passage au quotient.

Quand on est sur une base $S$, espace $\mathbb{C}$-analytique, la donnée a)
correspond : à un système local $\mathcal{H}$ sur $S$ --- faisceau en
$\mathbb{Z}$-modules loc. libres de rang $2$ sur $S$, d'où un fibré vectoriel
$\mathbb{R}$-analytique $\mathcal{H} \otimes_{\mathbb{Z}} \mathbb{R}$ sur $S$. La
structure b) correspond à une structure complexe \struck{sur} de ce dernier,
\uncertain{en disant que} fibré vectoriel \add{complexe} de rang $1$, de façon
que les sections locales de $\mathcal{H}$ deviennent des sections analytiques
complexes. Si \struck{\ill{}} $z_1, z_2$ sont deux sections de $\mathcal{H}$ formant

\page{143}
une base de $\mathcal{H} \otimes_{\mathbb{Z}} \mathbb{R}$ (\add{fibre par fibre}
sur $\mathbb{R}$), la donnée d'une telle structure complexe équivaut à la donnée de
\[
(3)\qquad \tau = z_2 / z_1
\]
une section du faisceau structural $\mathcal{O}_S$. (\emph{NB}
\struck{\ill{}}. Pour une structure complexe donnée sur
$\mathcal{H} \otimes_{\mathbb{Z}} \mathbb{R}$, $z_1$ et $z_2$ sont deux
sections\add{-base} des faisceau\supplied{x} inversible\supplied{s} \ill{}, et
$\tau$ est la section de $\mathcal{O}_S$ qui fait passer de l'une à l'autre :
\[
(3\,\mathrm{bis})\qquad z_2 = \tau z_1 .
\]
\marginal{Noter que $\forall s \in S$, $\tau(s) \in \mathbb{C} \setminus \mathbb{R}$}

Une structure \emph{de Torelli} sur une famille de courbes elliptiques relative
$E$ sur $S$ \uncertain{(fournissant} $\mathcal{H}$ sur $S$ $+$ structure complexe
sur $\mathcal{H} \otimes_{\mathbb{Z}} \mathbb{R}$) est la donnée d'un isomorphisme
\[
(4)\qquad t : \mathbb{Z}_S^2 \xrightarrow{\;\sim\;} \mathcal{H},
\]
\struck{\ill{}} \add{ce qui revient à} la donnée de deux sections $z_1, z_2$ de
$\mathcal{H}$, \struck{\uncertain{distinguées}} formant une base de $\mathcal{H}$
en ch. fibre. La catégorie des courbes elliptiques relatives sur $S$, avec
structure de Torelli, est « rigide » \struck{\ill{}} (pas d'automorphisme
$\neq \mathrm{id}$ d'un objet) et l'ens. des

\page{144}
classes d'isomorphie d'objets d'icelle est en corr. biunivoque canonique avec
l'ens. des applications holomorphes
\[
(5)\qquad \tau : S \longrightarrow \mathbb{C} \setminus \mathbb{R} .
\]
On dit qu'une structure de Torelli est « \uncertain{directe} » --- « stricte » si
\struck{\ill{}} pour $\forall s \in S$, la base $z_1, z_2$ de
$\mathcal{H}_s \otimes_{\mathbb{Z}} \mathbb{R}$ \struck{correspond} \add{définit}
l'orientation de cet espace correspondant à sa structure complexe, i.e.\ si (5)
\ill{} une application
\[
(6)\qquad \tau : S \longrightarrow \mathfrak{D}^+ = \{ z \in \mathbb{C} \mid \Im z > 0 \}.
\]
Donc $\mathfrak{D}^+$ apparaît comme l'espace analytique modulaire pour la
structure « courbe elliptique relative $+$ str. de Torelli stricte sur espace
analytique variable », la famille modulaire étant donnée par
\[
(7)\qquad \mathfrak{X}_{\mathfrak{D}^+} = (\mathfrak{D}^+ \times \mathbb{C}) / \mathbb{Z}^2
\]
où $\mathbb{Z}^2$ opère sur $\mathfrak{D}^+ \times \mathbb{C}$ par
\[
(8)\qquad T_{m,n}(\tau, z) = (\tau,\, z + m + n\tau), \qquad
\tau \in \mathfrak{D}^+,\ z \in \mathbb{C} .
\]
On voit que l'on a
\[
(8)\qquad (\mathfrak{X}_{\mathfrak{D}^+})_\tau \simeq \mathbb{C}/(\mathbb{Z} + \mathbb{Z}.\tau)
\simeq \mathbb{C}^* / \text{sous-groupe engendré par } k = \exp(2i\pi\tau)
\]
(\uncertain{on trouve bien} $|k| < 1$).
\note{le numéro (8) est porté deux fois sur la page. Devant $T_{m,n}$, un petit
groupe biffé.}

\page{145}
Sur l'ens. des \struck{structures} classes d'isomorphie de structures de Torelli
\uncertain{directes}, le groupe \struck{\ill{}}
\[
(9)\qquad \mathrm{SL}(2,\mathbb{Z}) = \left\{ \begin{pmatrix} a & b \\ c & d \end{pmatrix}
\;\middle|\; \begin{array}{l} a,b,c,d \in \mathbb{Z} \\ ad - bc = 1 \end{array} \right\}
\]
opère à droite, par
\[
t \longmapsto t \circ g ,
\]
on va expliciter cette opération, en tant qu'opération sur $\mathfrak{D}^+$, et
\ill{} sur $\mathfrak{X}_{\mathfrak{D}^+}$ (ce qui \ill{} une opération de
$\mathrm{SL}(2,\mathbb{Z})$ sur la situation modulaire).

On voit que : une structure de Torelli $(E, (z_1, z_2))$, $g$ amène la
structure $(E, (z'_1, z'_2))$, avec
\[
\begin{aligned}
z'_1 &= a z_1 + c z_2 \\
z'_2 &= b z_1 + d z_2
\end{aligned}
\]
d'où, si $\tau, \tau'$ sont les invariants correspondants ---
$\tau = z_2/z_1$, $\tau' = z'_2/z'_1$ ; la formule
\[
(10)\qquad \tau' = \frac{d\tau + b}{c\tau + a}
\qquad \left(\tau' = \tau \circ g,\ g = \begin{pmatrix} a & b \\ c & d \end{pmatrix}\right)
\]
Si on remplace l'opération à droite par l'opération à gauche par
\uncertain{$g$} \ill{} correspondante, \struck{\ill{}}
$t \longmapsto t \circ g^{-1}$, \ill{} $\tau \longmapsto \tau \circ g^{-1}$,
\uncertain{on trouve},
\note{le $c$ du dénominateur de (10) est écrit par-dessus une autre lettre.}


\page{146}
Au niveau de la famille modulaire, l'opération $t \longmapsto t \circ g$ se
calcule \ill{} qui correspond : \uncertain{explicitement}, si on a
\[
\begin{aligned}
E_\tau &= \mathbb{C}/(\mathbb{Z} + \mathbb{Z}\tau), \\
E_{\tau\circ g} &= \mathbb{C}/(\mathbb{Z} + \mathbb{Z}\tau')
\qquad \tau' = \tau \circ g = \frac{d\tau + b}{c\tau + a}
\end{aligned}
\]
on note que l'on a une application \add{$\mathbb{C}$}-linéaire
$\mathbb{C} \to \mathbb{C}$, qui transforme $\mathbb{Z} + \mathbb{Z}\tau$ en
$\mathbb{Z} + \mathbb{Z}\tau'$, qui est la multiplication par
$\lambda \bigl(= z_1/z'_1 = \dfrac{z_1}{a z_1 + c z_2} = \dfrac{1}{a + c\tau}\bigr)$,


\begin{tikzcd}
\mathbb{C} \arrow[r, "\sim"] \arrow[d, "2"'] & t_E \arrow[d, no head, "\Vert" description] \\
\mathbb{C} \arrow[r] & t_E
\end{tikzcd}

\note{schéma redessiné. Au-dessus de la flèche du haut :
$\lambda \mapsto \lambda z_1$ ; $1 \mapsto z_1$, $\tau \mapsto z_2$. Au-dessous de
celle du bas : $\lambda \mapsto \lambda z'_1$ ; $1 \mapsto z'_1$,
$\tau' \mapsto z'_2$. À gauche de la flèche verticale :
$\lambda \mapsto \lambda \frac{z_1}{z'_1}$ ; le « $2$ » qui l'étiquette est lu tel
quel. La flèche de droite est une égalité. En travers du schéma, de biais :
« \emph{NB} on a \uncertain{$1 = a\xi + c\xi\tau$, $\tau' = b\xi + d\xi\tau$} »,
la lettre notée $\xi$ n'étant pas sûre.}

et par passage au quotient il en résulte
\[
E_{\tau} \xrightarrow{\;\sim\;} E_{\tau\circ g} .
\]
Remplaçant $g = \begin{pmatrix} a & b \\ c & d \end{pmatrix}$ par
$\tau_\infty(g^{-1}) = \begin{pmatrix} d & b \\ c & a \end{pmatrix}$, on trouve
\[
E_\tau \xrightarrow{\;\sim\;} E_{\tau'}
\qquad \tau' = \tau \circ \tau_\infty(g^{-1}) = \frac{a\tau + b}{c\tau + d}
\]
\note{après $E_{\tau'}$, un « $= \ldots$ » est biffé.}
\add{réduit par} multiplication par $(c\tau + d)^{-1}$,
donc on fait opérer $g \in \mathrm{SL}(2,\mathbb{Z})$ \add{à gauche} sur
$\mathfrak{X}_{\mathfrak{D}^+} \simeq (\mathfrak{D}^+ \times \mathbb{C}) / \mathbb{Z}\times\mathbb{Z}$
via l'action
\[
(14)\qquad T_g.(\tau, z) = \Bigl(g.\tau,\ \frac{z}{c\tau + d}\Bigr)
\qquad \text{où } g = \begin{pmatrix} a & b \\ c & d \end{pmatrix},\
g.\tau = \frac{a\tau + b}{c\tau + d} .
\]
\note{devant $T_g$, un groupe biffé ; dans le second membre, un signe biffé
devant la fraction. Les formules (11) à (13) sont sur la p.~147. La matrice $\tau_\infty(g^{-1})$ est transcrite
telle qu'écrite.}


\page{147}
\note{cette page fait suite à la p.~145 (« on trouve ») ; la
p.~146, qui porte la formule (14), en est la suite. Les deux feuillets
sont rangés dans l'ordre inverse de celui du raisonnement.}
on trouve que
\[
g^{-1} = \begin{pmatrix} d & -b \\ -c & a \end{pmatrix}
\]
la formule
\[
(11)\qquad \tau \circ g^{-1} = \frac{a\tau - b}{-c\tau + d} .
\]
Introduisons l'\ill{}
\[
(12)\qquad
\begin{cases}
\tau_\infty = \begin{pmatrix} -1 & 0 \\ 0 & +1 \end{pmatrix} \in \mathrm{GL}(2,\mathbb{Z})
\quad \text{d'où} \\[6pt]
\tau_\infty \begin{pmatrix} A & B \\ C & D \end{pmatrix} \tau_\infty^{-1}
= \tau_\infty\Bigl(\begin{pmatrix} A & B \\ C & D \end{pmatrix}\Bigr)
= \begin{pmatrix} A & -B \\ -C & D \end{pmatrix} \\[6pt]
\tau_\infty\Bigl(\begin{pmatrix} a & b \\ c & d \end{pmatrix}^{-1}\Bigr)
= \begin{pmatrix} d & b \\ c & a \end{pmatrix}
\quad \text{« échange de } a \text{ et de } d \text{ »,}
\end{cases}
\]
on trouve
\[
(13)\qquad \tau \circ \tau_\infty(g^{-1}) = \frac{a\tau + b}{c\tau + d}
\qquad \text{où } g = \begin{pmatrix} a & b \\ c & d \end{pmatrix},\quad
\tau_\infty(g^{-1}) = \begin{pmatrix} d & b \\ c & a \end{pmatrix}
\]
Par la suite, quand on fait opérer $\mathrm{SL}(2,\mathbb{Z})$ sur
$\mathfrak{D}^+$, ce sera toujours (\uncertain{sauf} mention expresse) par la
formule (13) --- ce n'est donc pas l'action « modulaire » la plus évidente, mais
celle consacrée par l'usage, qui correspond \add{(\uncertain{à droite})} à la
précédente, modifiée via le homomorphisme
\[
\begin{pmatrix} a & b \\ c & d \end{pmatrix} = g \longmapsto
\tau_\infty(g^{-1}) = \begin{pmatrix} d & b \\ c & a \end{pmatrix} .
\]


\page{148}
de $\mathrm{SL}(2,\mathbb{Z})$ sur $\mathfrak{D}^+ \times \mathbb{C}$. On peut en
fait faire opérer $\underbrace{\mathbb{Z}^2 \cdot \mathrm{SL}(2,\mathbb{Z})}_{\text{produit semi-direct}}$
\struck{\ill{}} sur $\mathfrak{D}^+ \times \mathbb{C}$, par
\[
(15)\qquad T_{g,m,n}(\tau, z) = \Bigl(g\tau,\ \frac{z}{c\tau + d} + m + n\, g\tau\Bigr)
\]
et l'opération déduite de $\mathrm{SL}(2,\mathbb{Z})$ sur
$\mathfrak{X}_{\mathfrak{D}^+}$ se déduit de (15) en passant au quotient
\struck{\ill{}} dans $\mathfrak{D}^+ \times \mathbb{C}$ par le \struck{faux}
sous-groupe $\mathbb{Z}^2$ de $\mathbb{Z}^2 \cdot \mathrm{SL}(2,\mathbb{Z})$.

On peut énoncer dès maintenant

\emph{Théorème} La famille $\mathbb{C}$-analytique de courbes elliptiques,
\struck{définie} \add{modulaire} pour la rigidification de Torelli stricte, définie
par
\[
\begin{array}{c}
\mathfrak{X}_{\mathfrak{D}^+} = \mathfrak{D}^+ \times \mathbb{C} / \mathbb{Z}^2 \\
\Big\downarrow \scriptstyle (\tau, z) \mapsto \tau \\
\mathfrak{D}^+
\end{array}
\qquad
\begin{array}{l}
\mathbb{Z}^2 \text{ opérant par} \\
T_{m,n}(\tau, z) = (\tau, z + m + n\tau)
\end{array}
\]
est un revêtement universel de $M_{1,1}$, \struck{\uncertain{i.e.\ elle est
localement modulaire}, \ill{}} \uncertain{avec} groupe d'automorphismes sur
$M_{1,1}$ le groupe $\mathrm{SL}(2,\mathbb{Z})$, opérant par
\[
g\tau = \frac{a\tau + b}{c\tau + d} \qquad
g = \begin{pmatrix} a & b \\ c & d \end{pmatrix},\quad
\begin{array}{l} a,b,c,d \in \mathbb{Z} \\ ad - bc = 1 \end{array}
\]

L'opération précédente est aussi décrite, en termes du pb modulaire de Torelli, par


\page{149}
\[
g \cdot \tau = \tau \circ \tau_\infty(g^{-1})
\qquad \text{où}\quad
\begin{array}{l}
\tau_\infty = \begin{pmatrix} -1 & 0 \\ 0 & 1 \end{pmatrix} \in \mathrm{GL}(2,\mathbb{Z}) \\[4pt]
\tau_\infty(g^{-1}) = \begin{pmatrix} d & b \\ c & a \end{pmatrix}
\ \text{si } g = \begin{pmatrix} a & b \\ c & d \end{pmatrix}
\end{array}
\]

\emph{Corollaire 1} Le groupe fondamental $\mathcal{T}_{1,1}$ de $M_{1,1}$
relatif au revêtement universel précédent est canoniquement isomorphe à
$\mathrm{SL}(2,\mathbb{Z})$ (car les opérations précédentes \add{\uncertain{à gauche}}
de $\mathrm{SL}(2,\mathbb{Z})$ sur le revêtement, déduites de l'opération
\ill{} par le passage à \ill{}).

Les \ill{} sur les $\mathcal{T}_{g,\nu}$ donnent alors :

\emph{Corollaire 2} Soit $(X, o)$ une surface orientée compacte de genre $1$,
$\mathcal{T}$ le groupe des classes d'isotopie d'automorphismes de $(X, o)$,
$H = \pi_1(X, o) = H_1(X, \mathbb{Z})$, \struck{\uncertain{considérons}}
\add{alors} l'homomorphisme canonique
\[
(16)\qquad \mathcal{T} \longrightarrow \mathrm{Aut}^+_{\mathbb{Z}}(H)
\]
est un isomorphisme.

\emph{Corollaire 3} Soient $\pi$ un groupe \ill{} de type $(1,1)$,
$H = \pi_{\mathrm{ab}}$, $\mathcal{T}$ le groupe des automorphismes \uncertain{ext.}\ :
\ill{} de $\pi$ compatibles avec l'orientation, alors l'hom canonique
$\mathcal{T} \longrightarrow \mathrm{Aut}^+_{\mathbb{Z}}(H)$ est un iso.


\page{150}
Considérons maintenant des courbes elliptiques $E$, \add{\uncertain{munies}}
d'une « \emph{rigidification} différentielle », définie par une base
$\varphi$ de $t_E$, permettant donc d'identifier $t_E$ à $\mathbb{C}$. La donnée
d'une courbe elliptique \uncertain{munie} d'une rigidification différentielle
s'identifie donc : à la donnée d'un \emph{réseau} $H \subset \mathbb{C}$ de
$\mathbb{C}$. La donnée d'une courbe elliptique avec rigidification double Torelli
stricte $+$ différentielle, s'identifie donc à la donnée de
\[
(17)\qquad z_1, z_2 \in \mathbb{C}^* \quad \text{avec}\quad \tau = z_2/z_1 \in \mathfrak{D}^+
\]
i.e.\ \struck{la $\mathbb{R}$-base} $z_1, z_2$ formant \struck{\ill{}} une
$\mathbb{R}$-base de $\mathbb{C}$, qui est \emph{directe} pour l'orientation
canonique de $\mathbb{C}$. On fait opérer $\mathrm{SL}(2,\mathbb{Z})$ \uncertain{à
gauche} sur la variété analytique complexe
\[
(18)\qquad P\mathfrak{D}^+ = \{ (z_1, z_2) \mid z_1, z_2 \in \mathbb{C}^*,\ z_2/z_1 \in \mathfrak{D}^+ \}
\]
par
\[
(19)\qquad g(z_1, z_2) = (d z_1 + c z_2,\ b z_1 + a z_2)
\]
\marginal{c'est donc, via le passage (18) $g \mapsto \tau_\infty(g^{-1})$,
l'opération naturelle \ill{} \uncertain{au} pb modulaire \ill{}, qui
\ill{}}
\note{note marginale écrite de biais dans l'angle inférieur gauche, qui
recouvre en partie le numéro (19) ; ce numéro est lui-même écrit sur un (18).}
\uncertain{d'autre} famille modulaire bi-rigidifiée $\mathfrak{X}_{P\mathfrak{D}^+}$
peut s'écrire

\page{151}
\[
(20)\qquad \mathfrak{X}_{P\mathfrak{D}^+} \simeq P\mathfrak{D}^+ \times \mathbb{C} / \mathbb{Z}^2
\]
où $\mathbb{Z}^2$ opère sur $P\mathfrak{D}^+ \times \mathbb{C}$ par
\[
T_{m,n}(z_1, z_2, z) = (z_1, z_2, z + m z_1 + n z_2)
\]
Alors le groupe $\mathbb{C}^* \times \mathrm{SL}(2,\mathbb{Z})$ opère sur
$\mathfrak{X}_{P\mathfrak{D}^+} \to P\mathfrak{D}^+$, en faisant opérer
$\mathrm{SL}(2,\mathbb{Z})$ par
\[
(21)\qquad T_g(z_1, z_2, z) = (d z_1 + c z_2,\ b z_1 + a z_2,\ z)
\]
et $\mathbb{C}^*$ \struck{\ill{}} par
\[
(22)\qquad T_\lambda(z_1, z_2, z) = \Bigl(\frac{1}{\lambda} z_1,\ \frac{1}{\lambda} z_2,\ \frac{1}{\lambda} z\Bigr).
\]
\marginal{\uncertain{$\mathbb{C}^*$ opération des dilatations définie par} (22)
\ill{} \uncertain{de la base $\varphi$ de $t_E$ par $\lambda\varphi$} \ldots}
\note{note marginale écrite de biais dans la marge gauche, à la hauteur de (22) ;
une partie en est surchargée.}

Chacun des deux groupes \uncertain{commute} \struck{\ill{}} \uncertain{opérant}
librement sur $\mathfrak{X}_{P\mathfrak{D}^+}$ et \uncertain{même} sur
$P\mathfrak{D}^+$, mais l'opération du \uncertain{groupe} produit
$\mathbb{C}^* \times \mathrm{SL}(2,\mathbb{Z})$ \add{sur $P\mathfrak{D}^+$}
n'est pas libre, p.ex.\ l'élément
$(-1, -1)$ opère trivialement sur $P\mathfrak{D}^+$ (\uncertain{\ill{}} universel
des courbes elliptiques) \struck{\ill{}}
\struck{\ill{} mais \ill{} sur $\mathfrak{X}_{P\mathfrak{D}^+}$,
\ill{} $\mathfrak{X}_{P\mathfrak{D}^+}$ (\ill{} d'ordre $2$ !) \ill{}
opérations \ill{} $\mathbb{C}^*$ \ill{} et $\mathfrak{X}_{\mathfrak{D}^+}$ et le
quotient \ill{} la section nulle, \ill{} \ill{} $\mathrm{SL}(2,\mathbb{Z})$,
s'identifient \ill{} $M_{1,1}$ !}
\note{la fin de la page est biffée : plusieurs traits horizontaux et obliques
barrent un passage surchargé, où se lisent $\mathfrak{X}_{P\mathfrak{D}^+}$,
« d'ordre 2 ! », $\mathbb{C}^*$, $\mathfrak{X}_{\mathfrak{D}^+}$, « section
nulle », $\mathrm{SL}(2,\mathbb{Z})$ et, encadré, « $M_{1,1}$ ! ». Le passage
s'interrompt ; la page suivante est blanche.}


\page{153}
\note{la page s'ouvre sur un calcul de groupes, d'une autre plume, traversé de
longs traits obliques ; il est sans lien apparent avec ce qui précède et avec le
texte qui suit le trait horizontal.}
\[
1 \longrightarrow N \longrightarrow G \longrightarrow \mathfrak{S}_A \longrightarrow 1
\]
\[
N \simeq \bigl((\mathbb{F}_2)^A\bigr)' = \operatorname{Ker}\bigl(\mathbb{F}_2^A \to \mathbb{F}_2\bigr),
\qquad (x_i) \longmapsto \textstyle\sum x_i
\]
\note{au-dessus de $G$, une flèche courbe ramène $\mathfrak{S}_A$ vers un
$\mathfrak{S}^0 \subset G$ (lecture incertaine) ; sous $\mathfrak{S}_A$, un
signe d'inclusion et un second $\mathfrak{S}$ dont l'indice, surchargé, n'est
pas lu. L'indice de $\mathfrak{S}_A$ dans la suite exacte est lui aussi
surchargé.}
\[
\begin{aligned}
g &= u \cdot x = y u \qquad (u \in \mathfrak{S}_0,\ x \in N), &\quad
ux &= yu,\ \ y = u x u^{-1} = u(x) \\
g' &= u' \cdot x' = y' u' &\quad g'g &= y' \underline{u' y u'^{-1}}\, \underline{u' u} \\
g' g &= u' x' u x = (u' u)\bigl(\underbrace{u^{-1} x' u}_{u^{-1}(x')}\bigr) x = U'' X'' ,
&\quad U'' &= u'u,\ \ X'' = u^{-1}(x')\, x \\
&= \bigl(y' u'(y)\bigr)\, u' u
\end{aligned}
\]

\note{un trait horizontal sépare ce calcul de ce qui suit.}

L'extension quadratique de $\mathbb{Q}(t)$ est induite par l'extension
quadratique
\[
B = \mathbb{Z}[t]\bigl[S\bigr]/\bigl(S^2 - Q(t)\bigr)
\]
de $A = \mathbb{Z}[T]$. On voit que pour $\forall\, p \in P \setminus \{2\}$,
localisant $A$ en $pA$ d'où $A_{pA} = \mathbb{Z}_{p\mathbb{Z}}\{t\} \overset{\mathrm{df}}{=} A(p)$,
\struck{\ill{}} $B(p) \overset{\mathrm{df}}{=} B \otimes_A A(p)$ est normal (\uncertain{car}
$p$ divise $Q$, \ill{}). Si $p = 2$, alors $B(2)$ est normal si $2 \mid Q$. Si
$2 \nmid Q$, alors il est normal aussi si $Q \bmod 2 \in \mathbb{F}_2[t]$ n'est
pas un carré dans $\mathbb{F}_2[t]$ (donc pas non plus dans $\mathbb{F}_2(t)$).
\struck{Si} Il le reste \uncertain{si} \ill{} il est un carré dans
$\mathbb{F}_2[t]$, mais pas dans $\mathbb{Z}/4\mathbb{Z}[t]$ --- i.e.\ un
certain invariant dans $\mathbb{Z}/2\mathbb{Z}[t]$ n'est pas nul. Donc le seul cas
où $B(2)$ ne soit normal est celui où $2 \nmid Q$ \struck{\ill{}} $Q(t)$ est un
carré mod $4$ (p.ex.\ $Q(t) = t^2 - 4 = (t+2)(t-2)$ ;
\uncertain{ou} $Q(t) = 4t + 1$, $\varepsilon = +1$ ; $Q(t) = 4t - 1$,
$\varepsilon = -1$).

Si on excepte ce cas --- on voit que $B$ est normal en
\uncertain{codimension} $1$ --- i.e.\ \add{$\operatorname{Spec} B$} qui en ses
\uncertain{fermés} \add{\uncertain{pts non}} \struck{\uncertain{irréductibles}}
normaux et ses pts non fermés \uncertain{(donc isolés)}. Donc à un
\struck{\ill{}} ensemble fini près de $\operatorname{Spec} A = (\operatorname{Spec}
\mathbb{Z}/\ldots$, $B$ représente la normalisée de $\mathbb{E}'_{\mathbb{Z}}$
dans l'extension quadratique envisagée --- et ceci reste encore le cas si on
remplace $\mathbb{Z}$ par $\mathbb{Z}[1/2]$.
\note{le passage sous le trait est d'une lecture difficile ; les lettres
$A$, $B$, $Q$, $t$ sont sûres, les liaisons verbales beaucoup moins. Le texte
n'a pas de début sur ce feuillet (« L'extension quadratique ») ni de suite
visible sur la page 154. La fin de la ligne « $\operatorname{Spec} A = \ldots$ »
est coupée au bord de la page.}


\page{154}
\emph{Normalité des extensions quadratiques} $A[S]/(S^2 - a) = B$

On suppose d'abord $A$ \uncertain{anneau} \uncertain{régulier} \add{\uncertain{local}},
soit $p$ la car.\ résiduelle. \marginal{\emph{NB} Si $p \neq 2$, alors $B$
est local ssi $a \in \mathfrak{m}$ \uncertain{ou} ($a \notin \mathfrak{m}$ et
$\dot a \notin k^2$).}
$p \neq 2$, alors $B$ normal ssi $a \notin \mathfrak{m}^2$ --- et cette condition
est nécessaire en tous cas (y compris $p = 2$). Si $p = 2$, distinguons deux cas :

\begin{enumerate}
\item[a)] $a \in \mathfrak{m}$ \quad alors $B$ normal ssi $a \notin \mathfrak{m}^2$
\item[b)] $a \notin \mathfrak{m}$ i.e.\ $a$ unité. Soit $\dot a$ son image dans
$k = k(A) = A/\mathfrak{m}$.
\end{enumerate}
\begin{enumerate}
\item[$1^\circ$)] $\dot a \notin k^2$, alors $B$ est normal (\uncertain{c'est un
anneau} \ill{} \uncertain{de corps résiduel} $k[\sqrt{\dot a}\,]$)
\item[$2^\circ$)] $\dot a \in k^2$, soit $\dot a = \dot\alpha^2$, où $\alpha \in A$
déterminé mod $\mathfrak{m}$ --- donc $a = \alpha^2 - b$, avec
$b \in \mathfrak{m}$ \add{$(= \alpha^2 - a)$}. Posons $T = S - \alpha$,
alors
\[
S^2 - a = (T + \alpha)^2 - a = T^2 + 2\alpha T + b
\]
\[
B \simeq A[T]/(T^2 + 2\alpha T + b), \qquad
\hat B \simeq \underbrace{A[[T]]}_{C}/(T^2 + 2\alpha T + b)
\]
$T^2$, $2\alpha T$ sont dans $\mathfrak{r}(C)^2$, donc
$T^2 + 2\alpha T + b \notin \mathfrak{r}(C)^2$ ssi $b \notin \mathfrak{r}(C)^2$
i.e.\ $b \notin \mathfrak{m}^2$, i.e.\ $\tilde b \neq 0$ dans
$\mathfrak{m}/\mathfrak{m}^2$.
\end{enumerate}
\marginal{\emph{NB} $\tilde b \in \mathfrak{m}/\mathfrak{m}^2$ ne dépend
que du choix \uncertain{de $\alpha$} \ill{}}
\marginal{\emph{NB} $2 \in \mathfrak{m} \subset \mathfrak{r}(C)$,
$T \in \mathfrak{r}(C)$}

Donc les cas de normalité sont les suivants :
$a \in \mathfrak{m} \setminus \mathfrak{m}^2$, \emph{ou} $a$ de la forme
$\alpha^2 - b$ avec $\alpha$ unité, $b \in \mathfrak{m} \setminus \mathfrak{m}^2$,
\emph{ou} $a \notin \mathfrak{m}$, $\dot a \notin k^2$.
\marginal{\uncertain{Dans les cas normaux} : \uncertain{$a \in \mathfrak{m}$ ou}
$a \bmod \mathfrak{m}^2$ \uncertain{non carré} d'un \uncertain{élément} de
$A/\mathfrak{m}^2$, \ill{}}
\note{les notes marginales de cette page sont écrites de biais dans la marge
gauche ; la dernière est barrée de plusieurs traits obliques.}

\emph{Application} aux rev.\ quadratiques de $\mathbb{P}^1_{\mathbb{Z}}$
(génériquement \struck{normales} \add{\uncertain{cas de} $\mathbb{E}^1_{\mathbb{Z}}$})
\struck{\ill{}} au-dessus de $\operatorname{Spec} \mathbb{Z}$.

Le revêtement \struck{au-dessus du} fibre \struck{\ill{}} quadratique du pt
générique, correspond à une ext.\ quadratique de $\mathbb{Q}(t)$, donc décrit par
un élément de $\mathbb{Q}(t)^*/\mathbb{Q}(t)^{*2}$. Or $\mathbb{Q}(t)^*/\{\pm 1\}$
\uncertain{a une} base formée des éléments suivants
\begin{enumerate}
\item[a)] Les nbs premiers $p \in P$ \quad ($P = \{2, 3, 5, 7, \ldots\}$)
\item[b)] Les polynômes $P = a_0 t^n + a_1 t^{n-1} + \cdots + a_n \in \mathbb{Z}[t]$,
avec : $a_0 > 0$, les $a_i$ premiers entre eux (i.e.\ $\forall p \in P$, $P \bmod p$
$\in \mathbb{Z}_p[t]$ est $\neq 0$), $P$ premier dans $\mathbb{Q}[T]$.
\end{enumerate}
Ainsi un élément de $\mathbb{Q}(t)^*/\mathbb{Q}(t)^{*2}$ est représenté de façon
unique par un produit
\[
(1)\qquad P = \underset{\varepsilon}{\pm}\, \underbrace{p_1 p_2 \cdots p_n}_{q}\,
\underbrace{P_1 \cdots P_\ell}_{Q}
\]
où les $p_i \in P$ sont distincts, et les $P_j$ comme dans b) \uncertain{distincts}.
Cela signifie aussi que
\[
P = \pm q\, Q \qquad
\begin{cases}
q \in \mathbb{N}^* \text{ « quadratfrei »} \\
Q \in \mathbb{Z}[t] \text{ à coeff.\ dominant } > 0 \text{, « quadratfrei » dans } \mathbb{Q}[t] \\
\quad \text{à coeff.\ premiers entre eux.}
\end{cases}
\]
\note{« quadratfrei » est écrit en allemand. La lettre $P$ désigne à la fois
l'ensemble des nombres premiers et le polynôme de (1) ; on l'a laissée telle
quelle. Le $\mathbb{Z}_p[t]$ est sans doute $\mathbb{F}_p[t]$ ; on transcrit ce
qui est écrit.}

\page{155}
Considérons d'abord \uncertain{pour} $B \otimes \mathbb{F}_p$ sur
$A \otimes \mathbb{F}_p = \mathbb{F}_p[T]$, donné par l'équation
$S^2 - P(T) = 0$ \struck{\ill{}} sur $\mathbb{F}_p$
(\uncertain{supposons} $p \neq 2$). Cette extension est
\uncertain{génériquement} \struck{\ill{}} \add{étale} \add{\uncertain{ssi}}
$P(T) \not\equiv 0 \ (p)$. Cela montre déjà que \struck{\ill{}} l'on doit avoir,
dans (1), $q \in \{1, 2\}$, et dans le cas où $B(p)$ est \uncertain{normal}
\uncertain{(p.ex.\ si $q = 2$)} on vent donc que $B \otimes \mathbb{F}_2$ soit
\uncertain{génér.}\ étale sur $A \otimes \mathbb{F}_2$, on voit que l'on doit avoir
$q = 1$. Mais l'équation $S^2 - P(T)\ \bigl(= S^2 - \varepsilon Q(T)\bigr) = 0$ en
car.\ $2$ n'est jamais génériquement étale ! Donc \struck{il faut} pour
\uncertain{avoir} qqne chose génér.\ étale en \emph{tout cas}, il
\emph{faut} que $B(p)$ soit non normal de \uncertain{car.\ 2}, que
\[
P = \varepsilon\, Q \qquad \bigl(\varepsilon \in \{\pm 1\}\bigr)
\]
soit un carré mod $4$. \emph{NB} Pour $Q$ fixé, ceci fixe le
\emph{signe} $\varepsilon$ \uncertain{sans ambiguïté}, car $-1$ n'est pas un
carré mod $4$ !). Cette condition est aussi suffisante \ldots
\note{la ligne sous $P = \varepsilon Q$ porte un groupe fortement raturé ; on y
devine $\{\pm 1\}$, qui est transcrit.}

\emph{Exemple} (\uncertain{Légendre} modulaires des courbes elliptiques).
Soit $Q(T) = T(T - \lambda)$, avec $\lambda \in \mathbb{Z} \setminus \{0\}$
(dans le cas qui m'intéresse, $\lambda = 2^\alpha 3^\beta$ l'invariant modulaire
exceptionnel). On a
\[
Q(T) \equiv T^2 \ (2) \quad \text{si } \lambda \equiv 0 \ (4).
\]
On \uncertain{trouve} un revêtement étale de
\[
\mathbb{E}^1_{\mathbb{Z}} \setminus \{0, \lambda\} = \operatorname{Spec} \mathbb{Z}[T]_{T(T-\lambda)}
\]
\marginal{\ldots\ \uncertain{faut} $T' = 1/T$}
\uncertain{au-dessus de} $\infty$ de
\[
\mathbb{P}^1_{\mathbb{Z}} \setminus \{0, \lambda\} = \operatorname{Spec} \mathbb{Z}[T']_{\lambda T' - 1}
\]
\uncertain{on considère} $\mathbb{Z}[T']_{(\lambda T' - 1)}$, dans
$\mathbb{Z}[T']_{(\lambda T' - 1)}\bigl[S\bigr]/\bigl(S^2 - (\lambda T' - 1)\bigr)$.
\note{la page s'arrête sur cette formule. Dans la ligne $Q(T) \equiv T^2\ (2)$,
le terme de droite est surchargé (un $T^2$ ou $T^3$ écrit sur autre chose).}

\page{156}
\note{calculs griffonnés dans les blancs d'une convocation administrative
imprimée, tamponnée « 9 JUIN 1981 » et « 16 JUIN 1981 » ; seuls les calculs sont
transcrits. Ils reprennent la forme $T^2 + 2\alpha T + b$ de la
p.~154.}
\[
T^2 + 2\alpha T + b
\]
\[
x^2 + \alpha x y + \beta y^2, \qquad (x + \beta y)\ \ldots
\]
\[
\mathbb{F}_2[T, \pi] \qquad T^2 + \alpha \pi T + \beta \pi^2, \qquad
\alpha^2 - 4\beta = \alpha^2
\]
\[
4T - 1
\]
\note{dans $x^2 + \alpha x y + \beta y^2$ et $(x + \beta y)$, le $\beta$ est
écrit par-dessus un autre signe ; sous $T^2 + 2\alpha T + b$ figure un second
$T^2 +$ inachevé. Le $4$ de $4T - 1$ est surchargé.}

\end{document}
